\section{แนวคิดและหลักการที่เกี่ยวข้อง}
\subsection{สถาปัตยกรรมแบบแยกส่วน}
ระบบแยกส่วนติดต่อผู้ใช้ บริการประมวลผล และฐานข้อมูลออกจากกัน ส่วนหน้าเว็บรับข้อมูลและนำเสนอผล ขณะที่เซิร์ฟเวอร์รับผิดชอบการตรวจสิทธิ์และกฎการทำงาน โครงสร้างโครงการแบ่งเป็นแอปพลิเคชันเว็บและบริการ NestJS ภายใน monorepo เดียว ช่วยแยกความรับผิดชอบของแต่ละส่วนและจัดการการพัฒนาร่วมกัน \cite{repo}
\subsection{สิทธิ์และกระบวนการอนุมัติ}
การตรวจเฉพาะบทบาทไม่เพียงพอต่อการเข้าถึงใบลา ผู้ใช้ต้องมีความสัมพันธ์กับรายการนั้นด้วย เช่น เป็นเจ้าของคำขอหรือเป็นผู้อนุมัติที่ได้รับมอบหมาย ระบบจึงต้องตรวจทั้งตัวตน บทบาท และขอบเขตข้อมูลที่เซิร์ฟเวอร์ ส่วนการอนุมัติหลายลำดับต้องแยกสถานะของคำขอออกจากสถานะของแต่ละขั้น เพื่อระบุผู้ที่พิจารณาได้ในขณะนั้น
\subsection{ข้อมูลธุรกรรมและข้อมูลรายงาน}
ข้อมูลธุรกรรมใช้บันทึกคำขอและผลตัดสิน ส่วนข้อมูลรายงานใช้รวมยอดตามมิติ เช่น นิสิต ประเภทลา ภาคการศึกษา และรายวิชา การกำหนดความหมายของหนึ่งแถวให้ชัดเจนช่วยป้องกันการนับใบลาซ้ำเมื่อหนึ่งคำขอเกี่ยวข้องกับหลายวิชา คลังข้อมูลรุ่นปัจจุบันของโครงการแยกข้อมูลระดับคำขอ รายวิชา และขั้นอนุมัติ พร้อมบันทึกสถานะการโหลดข้อมูล \cite{warehouse}

\subsection{ความถูกต้องของข้อมูลข้ามระบบ}
ข้อมูลที่ใช้ยื่นคำขอและข้อมูลที่ใช้วิเคราะห์มีช่วงเวลาอ้างอิงต่างกัน รายชื่อผู้สอนหรือที่ปรึกษาปัจจุบันใช้ตรวจความสัมพันธ์ในการให้บริการ ขณะที่บริบทของคำขอที่ยื่นแล้วควรอธิบายด้วยหลักฐาน ณ เวลายื่น การเปลี่ยนข้อมูลอ้างอิงในภายหลังจึงไม่ควรทำให้ประวัติเดิมเปลี่ยนความหมายโดยไม่บันทึกที่มา แนวคิดนี้เป็นเหตุผลของการเก็บ snapshot ร่วมกับฐานข้อมูลธุรกรรมและการจัดทำ mapping สำหรับคลังข้อมูล
ความถูกต้องยังขึ้นกับการแยกรหัสแต่ละชนิด รหัสบัญชี รหัสบุคลากร รหัสนิสิต และคีย์ฐานข้อมูลอาจอ้างถึงบุคคลเดียวกันแต่ใช้แทนกันโดยตรงไม่ได้ การเชื่อมข้อมูลจึงต้องมีความสัมพันธ์ที่ตรวจสอบได้ หากพบรหัสกำกวมหรือข้อมูลที่ไม่ทราบ ควรรักษาสถานะนั้นไว้เพื่อแก้ต้นทาง แทนการจับคู่จากชื่อหรือเติมค่าที่ดูเป็นไปได้
